\section{Introduction}\label{sec:intro}

An HTTP server must decide which code answers each request. Its router takes the request's method
and path and returns the handler of the route that matches, with the values that the route's
parameters capture. The lookup runs once per request, beside parsing and input and output. A
server that answers many small requests on few cores pays for it on each one.

Web frameworks in C, C++, Rust and Go each ship a router of their own, and the routers differ in
structure. Some compress the routes into a radix
tree~\cite{schmidt2019httprouter,gin2026gin,chi2026chi,ahmed2026matchit,viz2025pathtree}. Others
build a trie of whole path segments~\cite{pistache2024pistache,uwebsockets2026uwebsockets,berry2026glaze}
or a prefix trie that uses regular expressions for patterns~\cite{r32026r3}. Crow builds a character
trie and merges chains of single children~\cite{crow2026crow}. Still others try a list of patterns in
order~\cite{oatpp2024oatpp,hirose2026httplib,actix2026router}, or match a regular expression per
parameter route~\cite{drogon2026drogon}. They also differ in what a route means: which of two
matching routes wins, whether a trailing slash matters, and what a catch-all captures. The
peer-reviewed literature on router speed is thin. We found two papers on routing with
tries~\cite{zhang2015urlrouting,dsouza2018htu}. We also found benchmark suites kept beside
routers~\cite{schmidt2020gohttproutingbench,ahmed2026matchit,sworn2026wayfind,vm0012024gateways}.

The author's earlier, co-authored work introduced RegexMatcher~\cite{stankov2024regex}. It matches
one string against many regular expressions at once, and web servers whose routes are regular
expressions motivated it. A route written as a regular expression loses what a router needs,
however. Precedence then follows the order of registration, a typed parameter is only a character
class, and every lookup runs an automaton over the whole path. This paper reworks RegexMatcher into
a route matcher, RegexMatcher~v2, and keeps the regular-expression engine, v1, beside it in the same
library.

The paper makes four contributions:
\begin{itemize}
\item A survey of the route matchers of sixteen routers written in C, C++, Rust and Go. It states
how each one stores its routes and what its lookup returns. It then measures all of them on the
same tables, with every answer checked, and runs eight probes of the path rules on each.
\item RegexMatcher~v2, a route matcher in C++23. A method whose routes are all literal gets an
exact-match table, and every other method a trie of path segments. The tables live in flat arrays
and a lookup never allocates. Patterns, handler signatures and duplicate routes are checked when the program compiles.
\item A comparison under seven hypotheses that were pre-specified in the paper's repository before
the run, at commit \PreSpecCommit{} of its earlier history (Section~\ref{sec:prespec}). It covers lookup time, work per lookup, allocation, build
time and memory, the compile-time checks, an end-to-end server, and a table built at compile time.
\item The results as the pre-specified rules decide them, losses included.
\end{itemize}

Lookup speed is the main result. v2 was not slower than the fastest competitor in \HOnePassed{}
of \HOneCells{} cells, and faster in \HOneSuperior{}. Every such claim holds for one measured setting
only: one core of an \HostLCpu{} (\HostLArch) and \RunCompiler. Every arm is compiled for that
processor's own instruction set, and the competitors run in the configurations of
Appendix~\ref{app:config}. The compile-time checks come second.

Section~\ref{sec:related} surveys related work and the competing routers.
Section~\ref{sec:design} describes v2. Section~\ref{sec:methods} gives the pre-specified design of
the measurement. Section~\ref{sec:results} reports each hypothesis and the exploratory runs.
Section~\ref{sec:discussion} discusses the losses and the threats to validity.
Section~\ref{sec:conclusions} concludes.
